\documentclass[book.tex]{subfiles}

\begin{document}
\chapter{Case Studies of Neural Potentials}

\label{chap:case_studies}
\noindent\rule{11cm}{0.4pt} \\
\textbf{Prerequisites:}  \autoref{chap:molecular_hamiltonian},   \\
\textbf{Difficulty Level:} ** \\
\noindent\rule{11cm}{0.4pt}
\vspace{5mm}

"""
The session with begin with a discussion on how to enable material-property predictions by combining advanced statistical mechanics with data-driven machine learning interatomic potentials. Two example applications will be covered using first-principles thermodynamic descriptions. Then, the talk will focus on why machine learning potentials work at all, and in particular, the locality argument. We will see that a machine-learning potential trained on liquid water alone can predict the properties of diverse ice phases because all the local environments characterizing the ice phases are found in liquid water.
The code walkthrough will use the Python code: https://github.com/BingqingCheng/ASAP
"""

References \cite{cheng2019ab}, \cite{cheng2020evidence}

Bingqing Cheng's two papers

- Water phase diagram
- Water quantum properties

\url{https://scholar.google.ch/citations?user=s5ZqEskAAAAJ&hl=en} 

\end{document}